\documentclass[12pt]{article}
\usepackage[margin=1in]{geometry}
\usepackage{setspace}
\usepackage{amsmath,amssymb}
\usepackage{enumitem}
\usepackage{hyperref}
\usepackage{xcolor}
\setstretch{1.2}

\title{Memo: How to Incorporate the Polarization-Measurement Papers into the Current Draft}
\author{}
\date{\today}

\begin{document}
\maketitle

\section*{Bottom line}
All three papers are relevant, but they are relevant in \emph{different ways}. The strongest immediate use is not to replace the paper's current empirical design, but to make the draft's measurement strategy more disciplined, more transparent about what each proxy captures, and more careful about what the current data can and cannot identify.

The current draft already uses three empirical proxies: (i) a state-year electoral polarization measure based on the Esteban--Ray index applied to House election returns, (ii) a state-level and national DW-NOMINATE ideological-gap measure, and (iii) a supplementary national ANES affective-polarization series interacted with state exposure. In the current version, this architecture is workable, but the theoretical and measurement discussion is still too ad hoc. The three uploaded papers can fix that.

My assessment is:
\begin{enumerate}[label=\arabic*.]
    \item \textbf{Wang and Tsui is relevant mainly as a measurement-discipline paper.} It is useful for clarifying that polarization is not the same thing as inequality, for distinguishing \emph{increased spread} from \emph{increased bipolarity}, and for explaining why different polarization indices embody different ethical choices rather than being mechanically interchangeable.
    \item \textbf{Esteban, Grad\'{\i}n, and Ray (2007) is highly relevant as a bridge paper.} It explicitly shows how the original ER framework can be extended to densities, how clustering creates approximation error, and how Wolfson can be derived as a special case in a unified framework.
    \item \textbf{Duclos, Esteban, and Ray (2004) is conceptually the most important.} It gives the cleanest statement of the identity--alienation logic, shows why polarization is distinct from inequality, and provides the most rigorous density-based formulation and estimation framework. However, its methodology is not directly plug-and-play for the current event-study design unless you are willing to build a separate density-estimation project.
\end{enumerate}

So the right move is \textbf{not} to rebuild the paper around continuous-distribution polarization measurement. The right move is to use these papers to improve: (a) the conceptual framing of what the paper's proxies do, (b) the rationale for having multiple proxies, and (c) the interpretation of robustness tests.

\section*{I. What the three papers contribute}

\subsection*{1. Wang and Tsui (2000)}
Wang and Tsui start from the Foster--Wolfson ordering and emphasize that polarization has two distinct comparative-static components: \emph{increased spread} (mass moving away from the middle) and \emph{increased bipolarity} (mass within each side becoming more internally compact). They show that Lorenz dominance is not adequate for polarization because transfers that reduce inequality can increase polarization, and they derive equivalent formulations of the Foster--Wolfson ordering together with generalized classes of indices. In their setup, the core comparative-static language is median-centered and explicitly distinct from standard inequality orderings.

For the draft, the key takeaway is not the exact index formula. The key takeaway is that a polarization measure is doing conceptual work only if it is sensitive to both \emph{distance across groups} and \emph{cohesion within groups}. That gives you a principled way to discuss what the current electoral proxy does and does not capture.

\subsection*{2. Esteban, Grad\'{\i}n, and Ray (2007)}
Esteban, Grad\'{\i}n, and Ray explicitly describe the identity--alienation framework and then ask what to do when the data are not naturally grouped. Their answer is to choose a finite number of groups, endogenize the cutoffs by minimizing within-group dispersion, and define \emph{extended polarization} as ER polarization on the grouped representation minus an error term that captures information lost from clustering:
\begin{equation*}
P(f;\alpha,\beta)=ER(\alpha,\rho)-\beta\,\varepsilon(f,\rho).
\end{equation*}
They further show that Wolfson is a special case in the bipolar case.

This paper is directly useful for the current draft because it gives you a formal language for saying that any grouped empirical proxy is a \emph{representation} of a latent polarization concept, not the concept itself. That is exactly the problem your current draft faces: the state-year electoral proxy is a grouped political-outcome proxy for a deeper latent social environment.

\subsection*{3. Duclos, Esteban, and Ray (2004)}
Duclos, Esteban, and Ray is the strongest conceptual anchor. It argues that polarization differs from inequality because alienation alone is not enough; effective antagonism requires both \emph{alienation} and \emph{identification}. They characterize a density-based polarization class under an identity--alienation structure and show that the resulting measure depends on a polarization-sensitivity parameter $\alpha$. They also develop kernel-based estimators and distribution-free inference for the density case.

For your draft, the most important use of this paper is conceptual rather than mechanical. It gives you the cleanest possible justification for treating polarization as something more than issue distance. In your setting, heterogeneous interpretation of auditor disclosures can be framed as arising because polarization changes both perceived distance from out-groups and identification with in-groups or institutions.

\section*{II. How these papers map onto the current draft}

\subsection*{1. The current draft already has a workable empirical hierarchy}
Right now the draft uses:
\begin{enumerate}[label=\alph*.]
    \item a baseline state-year ER-style electoral measure,
    \item a DW-NOMINATE ideological-gap robustness measure,
    \item a supplementary ANES affective-polarization series interacted with state exposure.
\end{enumerate}
This is a sensible hierarchy for the current project. The problem is not that the draft lacks measures. The problem is that the draft does not yet explain \emph{why} these measures belong together.

The three papers should therefore be used to organize the measurement section around one sentence:

\begin{quote}
The model's polarization object is latent; the empirical proxies are complementary measurements of different margins of that object rather than interchangeable ``politics variables.''
\end{quote}

\subsection*{2. What each current proxy should be said to measure}
The draft should explicitly assign each empirical proxy a role:
\begin{itemize}
    \item \textbf{ER election-return measure:} a reduced-form proxy for \emph{mass ideological sorting and electoral competitiveness} in the local political environment.
    \item \textbf{DW-NOMINATE gap:} a proxy for \emph{elite ideological polarization}. This is conceptually closer to issue-position divergence than to partisan animus.
    \item \textbf{ANES feeling-thermometer series:} a proxy for \emph{national affective polarization}. This is the closest current measure to distrust, animus, and social distance.
\end{itemize}

Once stated this way, the measurement strategy becomes coherent. ER and DW-NOMINATE are \emph{not} two ways of measuring the same thing; one is electoral and local, the other is elite and ideological. The ANES series is not a substitute for either of them; it is a separate affective validation.

\section*{III. Specific drafting changes recommended}

\subsection*{A. Introduction}
The introduction should stop presenting the measurement strategy as if the empirical proxies directly observe investor beliefs. Instead, add a short paragraph after the contribution paragraph along the following lines:

\begin{quote}
A central empirical challenge is that the model's latent polarization object is not directly observable at the investor level. Following the polarization-measurement literature, we therefore use complementary proxies that capture distinct margins of polarization: local electoral polarization, elite ideological polarization, and national affective polarization. We interpret these as imperfect but informative measurements of the political environment in which auditor disclosures are received.
\end{quote}

This paragraph should cite Duclos--Esteban--Ray for the latent identity--alienation concept, Wang--Tsui for the fact that polarization is not equivalent to inequality, and Esteban--Grad\'{\i}n--Ray for the point that grouped empirical measures are representations that necessarily abstract from within-group dispersion.

\subsection*{B. Theory section}
The theory currently speaks broadly about heterogeneous priors, institutional trust, and perceived precision. That is fine, but it should explicitly separate \emph{the latent object} from \emph{the measured object}. Add a paragraph such as:

\begin{quote}
The theory treats polarization as a latent force that increases divergence in interpretation through identity, alienation, trust, and perceived signal precision. The empirical measures below do not observe this latent object directly. Instead, they proxy for different margins of it: electoral sorting in the local political environment, elite ideological distance, and national affective animus.
\end{quote}

This is where Duclos--Esteban--Ray belongs most naturally. Their framework is the cleanest support for the idea that polarization is not merely issue distance but a combination of identification and alienation.

\subsection*{C. Measurement subsection}
The current measurement subsection should be rewritten into three numbered subparts.

\paragraph{(i) Why grouped proxies are unavoidable.}
Use Esteban--Grad\'{\i}n--Ray to state that empirical proxies often work with grouped or clustered representations of a latent polarization object, and that such representations inevitably lose within-group information. This justifies why your state-year ER proxy is a tractable reduced-form measure rather than a literal measure of investor-level polarization.

\paragraph{(ii) Why the baseline ER measure is still useful.}
Use Wang--Tsui here. Their paper provides language for saying that a useful polarization proxy should reflect both separation and clustering, not just dispersion. You should then candidly note that the district-election ER proxy captures electoral division and local sorting, but does not separately identify within-group cohesion in the same way a continuous-density approach would.

\paragraph{(iii) Why alternative measures are needed.}
Use both Wang--Tsui and Duclos--Esteban--Ray to justify multiple proxies. Wang--Tsui says different indices embody different judgments and are not interchangeable; Duclos--Esteban--Ray says the latent concept is richer than any one reduced-form statistic. Therefore, the purpose of DW-NOMINATE and ANES is not cosmetic robustness but triangulation.

\section*{IV. What should \emph{not} be done}

\subsection*{1. Do not rebuild the paper around continuous-density estimation}
The Duclos and Esteban--Grad\'{\i}n--Ray methodologies are excellent, but they are not immediate substitutes for the current design. To use them directly, you would need a density over the relevant underlying political or investor-level trait. The current data do not give you a credible firm-year density of investor ideology or animus. So the memo should recommend against claiming that you are \emph{estimating polarization in the DER sense}. You are not.

\subsection*{2. Do not oversell the ER election-return proxy as a structural measure}
The current draft risks treating the ER election-return measure as if it were a direct estimate of polarization in the underlying investor environment. The uploaded papers imply the opposite lesson: any grouped proxy is a reduced-form stand-in for a richer latent object. So the memo should recommend explicitly narrowing the claim.

\subsection*{3. Do not say the paper separately identifies ideological and affective polarization at the state-year level}
The current ANES affective measure is national and only becomes usable through interaction with a cross-sectional exposure term. That is still useful, but it is not the same thing as a state-year affective panel. The draft should not imply otherwise.

\section*{V. Concrete empirical recommendations inspired by the papers}

\subsection*{1. Reframe Table 5 as measurement triangulation}
Instead of presenting the robustness table as a generic list, relabel and discuss it as a \emph{measurement triangulation table}:
\begin{itemize}
    \item Columns (1)--(3): ER sensitivity to different $\alpha$ values.
    \item Columns (4)--(5): alternative ideological measures using DW-NOMINATE.
    \item Column (6): state fixed effects or clustering refinement.
    \item Column (7): the ANES affective series interacted with exposure.
\end{itemize}
The interpretive claim should be: similar signs across these measures strengthen the claim that the effect is tied to polarization broadly rather than to one specific formula.

\subsection*{2. Add one short ``limitations of measurement'' paragraph}
This paragraph should say three things:
\begin{enumerate}[label=\arabic*.]
    \item No single empirical measure fully captures the latent identity--alienation object.
    \item Electoral proxies, elite ideological proxies, and affective proxies each capture different margins.
    \item Convergent evidence across them is therefore more informative than reliance on one baseline statistic.
\end{enumerate}

\subsection*{3. Add a footnote on why the paper does not implement DER directly}
A brief footnote could say:
\begin{quote}
The continuous-density estimators in Duclos, Esteban, and Ray (2004) require a directly observed distribution of the relevant underlying trait. Because our setting does not provide a firm-year density of investor political identity, we use reduced-form political-environment proxies instead.
\end{quote}
This heads off the objection that you cited DER but ignored its estimator.

\section*{VI. Suggested paragraph-level edits for the draft}

\subsection*{Edit 1: Replace the opening sentence of the measurement section}
Current approach: the section dives immediately into the ER formula.

Recommended replacement:
\begin{quote}
The empirical challenge is that polarization is latent, while the available data provide only grouped or indirect proxies. Following the polarization-measurement literature, we therefore use a baseline local electoral proxy and then test whether the results are robust to conceptually distinct ideological and affective measures.
\end{quote}

\subsection*{Edit 2: Add a paragraph after the ER formula}
\begin{quote}
This baseline measure should be interpreted as a reduced-form proxy for the political environment surrounding the firm rather than as a direct estimate of investor-level polarization. In the terminology of Esteban, Grad\'{\i}n, and Ray (2007), it is a grouped empirical representation of a richer latent object and therefore necessarily abstracts from within-group heterogeneity.
\end{quote}

\subsection*{Edit 3: Add a paragraph before DW-NOMINATE}
\begin{quote}
Because different polarization indices emphasize different aspects of between-group separation and within-group cohesion, the baseline electoral measure cannot be treated as exhaustive. Wang and Tsui (2000) show that polarization orderings depend on whether a measure is sensitive to increased spread, increased bipolarity, or both. We therefore supplement the baseline with a measure of elite ideological distance from congressional roll-call voting.
\end{quote}

\subsection*{Edit 4: Add a paragraph before the affective measure}
\begin{quote}
The identity--alienation framework in Duclos, Esteban, and Ray (2004) suggests that issue distance is only one component of polarization. Affective antagonism is conceptually distinct from ideological divergence and is especially relevant for a mechanism operating through distrust in auditors and regulatory institutions. For that reason, we use a supplementary national ANES affective-polarization series interacted with local partisan exposure.
\end{quote}

\section*{VII. Recommended bibliography additions and where they belong}
\begin{enumerate}[label=\arabic*.]
    \item \textbf{Wang and Tsui (2000)} should appear in the measurement subsection, especially around the discussion of why polarization is not equivalent to inequality and why different indices are not redundant.
    \item \textbf{Esteban, Grad\'{\i}n, and Ray (2007)} should appear in the measurement subsection and possibly in a short methodological footnote on grouped proxies and approximation error.
    \item \textbf{Duclos, Esteban, and Ray (2004)} should appear in the theory section and in the measurement section, because it is the best bridge between concept and measurement.
\end{enumerate}

\section*{VIII. My final recommendation}
These papers are relevant enough that they should be incorporated explicitly, but not in a way that changes the paper's core design.

The right integration strategy is:
\begin{enumerate}[label=\arabic*.]
    \item Use \textbf{Duclos--Esteban--Ray} to anchor the conceptual definition of polarization.
    \item Use \textbf{Wang--Tsui} to justify why multiple ideological proxies are informative and why polarization is distinct from inequality or simple dispersion.
    \item Use \textbf{Esteban--Grad\'{\i}n--Ray} to justify grouped empirical proxies and to explain why the current baseline is a tractable reduced-form proxy rather than a structural estimate.
\end{enumerate}

That will make the measurement strategy look much more deliberate, methodologically literate, and internally consistent, while keeping the paper implementable for the current project.

\section*{References}
\begin{thebibliography}{9}
\bibitem{} Duclos, Jean-Yves, Joan Esteban, and Debraj Ray. 2004. ``Polarization: Concepts, Measurement, Estimation.'' \emph{Econometrica} 72: 1737--1772.
\bibitem{} Esteban, Joan, Carlos Grad\'{\i}n, and Debraj Ray. 2007. ``An Extension of a Measure of Polarization, with an Application to the Income Distribution of Five OECD Countries.'' \emph{Journal of Economic Inequality} 5: 1--19.
\bibitem{} Wang, You-Qiang, and Kai-Yuen Tsui. 2000. ``Polarization Orderings and New Classes of Polarization Indices.'' \emph{Journal of Public Economic Theory} 2: 349--363.
\end{thebibliography}

\end{document}
